% !TEX root = ../Main.tex
\chapter{提高篇}

本章在基础篇之上加入\textbf{两类难点}：(1) 求某个参数（棱段比、线段长）使某个\textbf{几何条件}成立；(2) \textbf{折叠}后的坐标重建。套路仍然是"建系 + 法向量"，但\textbf{参数化}要更细致。

\section{提高 1：菱形 + 竖立矩形的面面角}

\ex[菱形底 + 矩形侧面]{
    如图，菱形 $ABCD$ 的边长为 $4$，$\angle DAB = 60^\circ$。矩形 $BDFE$ 与底面 $ABCD$ 垂直，且面积为 $8$。
    \begin{enumerate}
        \item[(1)] 求证：$AC \perp BE$；
        \item[(2)] 求二面角 $E$-$AF$-$D$ 的余弦值。
    \end{enumerate}
}

\sol{
    \textbf{准备}：$\angle DAB = 60^\circ$、边长 $4$ $\Rightarrow \triangle ABD$ 为等边，$BD = 4$。对角线 $AC = 2 \cdot 4 \cos 30^\circ = 4 \sqrt 3$。矩形 $BDFE$ 面积 $= |BD| \cdot |BE| = 4 |BE| = 8 \Rightarrow |BE| = 2$。

    \textbf{建系}：取菱形对角线交点 $O$ 为原点，$x$ 沿 $OA \to OC$ 方向，$y$ 沿 $OB \to OD$ 方向，$z$ 竖直。
    \[
    A = (-2\sqrt 3, 0, 0),\ B = (0, -2, 0),\ C = (2\sqrt 3, 0, 0),\ D = (0, 2, 0).
    \]
    矩形 $BDFE$ 在 $yz$ 平面内（$\perp$ 底面且含 $BD$），$BE \perp BD$、$|BE| = 2$：
    \[
    E = (0, -2, 2),\ F = (0, 2, 2).
    \]

    \textbf{(1)} $\vec{AC} = (4\sqrt 3, 0, 0)$，$\vec{BE} = (0, 0, 2)$。$\vec{AC} \cdot \vec{BE} = 0 \Rightarrow AC \perp BE$。$\blacksquare$

    \textbf{(2)} 平面 $EAF$：$\vec{AE} = (2\sqrt 3, -2, 2)$，$\vec{AF} = (2\sqrt 3, 2, 2)$。
    \[
    \vec n_1 = \vec{AE} \times \vec{AF} = (-4 - 4,\ 4\sqrt 3 - 4\sqrt 3,\ 4\sqrt 3 + 4\sqrt 3) = (-8, 0, 8\sqrt 3).
    \]
    取简化 $\vec n_1 = (-1, 0, \sqrt 3)$，$|\vec n_1| = 2$。

    平面 $DAF$：$\vec{AD} = (2\sqrt 3, 2, 0)$，$\vec{AF} = (2\sqrt 3, 2, 2)$。
    \[
    \vec n_2 = \vec{AD} \times \vec{AF} = (4 - 0,\ 0 - 4\sqrt 3,\ 4\sqrt 3 - 4\sqrt 3) = (4, -4\sqrt 3, 0).
    \]
    取简化 $\vec n_2 = (1, -\sqrt 3, 0)$，$|\vec n_2| = 2$。

    \[
    \vec n_1 \cdot \vec n_2 = -1 + 0 + 0 = -1.
    \]
    \[
    \cos \theta = \frac{-1}{2 \cdot 2} = -\frac{1}{4}.
    \]
    故二面角 $E$-$AF$-$D$ 的余弦值为 $\boxed{-\dfrac{1}{4}}$（钝角）。
}

\section{提高 2：参数化求位置使面面角指定}

\ex[参数化的面面角问题]{
    四棱锥 $P$-$ABCD$ 中，$PA \perp$ 底面 $ABCD$，底面 $ABCD$ 为直角梯形，$AD \parallel BC$，$AD \perp CD$，$AD = CD = 2\sqrt 2$，$BC = 4 \sqrt 2$，$PA = 2$。$M$ 在 $PD$ 上，记 $\dfrac{PM}{PD} = t$。求 $t$，使二面角 $M$-$AC$-$D$ 的余弦值为 $\dfrac{\sqrt 2}{2}$。
}

\sol{
    \textbf{建系}：$A$ 为原点，$x$ 轴沿 $CD$ 方向（$A$ 到 $CD$ 的平行线），$y$ 轴沿 $AD$，$z$ 轴沿 $AP$。

    \[
    A = (0, 0, 0),\ D = (0, 2\sqrt 2, 0),\ C = (2\sqrt 2, 2\sqrt 2, 0),\ P = (0, 0, 2).
    \]
    $M = P + t(D - P) = (0,\ 2\sqrt 2\, t,\ 2 - 2t)$。

    \textbf{二面角 $M$-$AC$-$D$} 沿边 $AC$。$\vec{AC} = (2\sqrt 2, 2\sqrt 2, 0)$，$|\vec{AC}|^2 = 16$。

    $\vec{AM} = (0, 2\sqrt 2 t, 2 - 2t)$。$\vec{AM} \cdot \vec{AC} = 8 t$。投影系数 $= \tfrac{8t}{16} = \tfrac{t}{2}$。$\vec{AM}$ 垂直 $\vec{AC}$ 的分量：
    \[
    \vec u = \vec{AM} - \tfrac{t}{2} \vec{AC} = (-\sqrt 2\, t,\ \sqrt 2\, t,\ 2 - 2t).
    \]
    $|\vec u|^2 = 2t^2 + 2t^2 + (2-2t)^2 = 8t^2 - 8t + 4$。

    $\vec{AD} = (0, 2\sqrt 2, 0)$，$\vec{AD} \cdot \vec{AC} = 8$，投影系数 $= \tfrac{1}{2}$。
    \[
    \vec v = \vec{AD} - \tfrac{1}{2} \vec{AC} = (-\sqrt 2, \sqrt 2, 0),\quad |\vec v| = 2.
    \]
    $\vec u \cdot \vec v = 2t + 2t + 0 = 4t$。
    \[
    \cos \theta = \frac{\vec u \cdot \vec v}{|\vec u|\,|\vec v|} = \frac{4t}{2 \sqrt{8t^2 - 8t + 4}} = \frac{2t}{\sqrt{8t^2 - 8t + 4}}.
    \]
    设 $= \dfrac{\sqrt 2}{2}$ 且 $t \in (0, 1)$（$t > 0$ 保证 $\cos > 0$）：
    \[
    \frac{4 t^2}{8 t^2 - 8 t + 4} = \frac{1}{2} \Rightarrow 8 t^2 = 8 t^2 - 8 t + 4 \Rightarrow t = \tfrac{1}{2}.
    \]
    故 $t = \dfrac{1}{2}$，即 $M$ 是 $PD$ 的\textbf{中点}。
}

\section{提高 3：矩形沿 $EF$ 折叠}

\ex[折叠求距离]{
    矩形 $ABCD$ 中，$AB = 2$，$AD = 4$；$E, F$ 分别在 $AD, BC$ 上使 $DE = CF = 1$。沿 $EF$ 将矩形右端的 $CDEF$ 部分折起到 $C'D'EF$，使二面角 $D'$-$EF$-$A$ 为 $60^\circ$。求 $D'$ 到平面 $ABFE$ 的距离。
}

\sol{
    \textbf{建系}：折叠后保持底面 $ABFE$ 不动。取 $A = (0, 0, 0)$，$x$ 沿 $AB$，$y$ 沿 $AE$（即 $AD$）：
    \[
    A = (0, 0, 0),\ B = (2, 0, 0),\ F = (2, 3, 0),\ E = (0, 3, 0).
    \]
    （$AE = AD - DE = 4 - 1 = 3$；$F$ 为 $BC$ 上距 $C$ 为 $1$ 即距 $B$ 为 $3$ 之点。）

    折叠前 $D$ 在 $y$ 正方向、距 $E$ 为 $1$。折叠后 $D'$ 绕轴 $EF$（沿 $x$ 轴方向）旋转 $60^\circ$（由二面角决定）：
    \[
    \vec{ED'} = (0,\ \cos 60^\circ,\ \sin 60^\circ) = \left(0,\ \tfrac{1}{2},\ \tfrac{\sqrt 3}{2}\right),
    \]
    即 $D' = E + \vec{ED'} = \left(0,\ \tfrac{7}{2},\ \tfrac{\sqrt 3}{2}\right)$。

    底面 $ABFE$ 在 $xy$ 平面内，$D'$ 到底面的距离即 $D'$ 的 $z$ 坐标绝对值：
    \[
    d = \dfrac{\sqrt 3}{2}.
    \]

    故 $D'$ 到平面 $ABFE$ 的距离为 $\boxed{\dfrac{\sqrt 3}{2}}$。
}

\section{提高 4：异面直线角}

\ex[异面直线角的小综合]{
    在正方体 $ABCD$-$A_1B_1C_1D_1$ 中，棱长为 $1$，$M, N$ 分别是 $A_1B, C_1D$ 的中点。求异面直线 $MN$ 与 $AC_1$ 所成角的余弦值。
}

\sol{
    \textbf{建系}：$A = (0, 0, 0),\ B = (1, 0, 0),\ C = (1, 1, 0),\ D = (0, 1, 0),\ A_1 = (0, 0, 1),\ B_1 = (1, 0, 1),\ C_1 = (1, 1, 1),\ D_1 = (0, 1, 1)$。
    \[
    M = \tfrac{1}{2}(A_1 + B) = (\tfrac{1}{2}, 0, \tfrac{1}{2}),\quad N = \tfrac{1}{2}(C_1 + D) = (\tfrac{1}{2}, 1, \tfrac{1}{2}).
    \]
    \[
    \vec{MN} = (0, 1, 0),\ |\vec{MN}| = 1;\quad \vec{AC_1} = (1, 1, 1),\ |\vec{AC_1}| = \sqrt 3.
    \]
    \[
    \vec{MN} \cdot \vec{AC_1} = 1.
    \]
    \[
    \cos \angle = \frac{|1|}{1 \cdot \sqrt 3} = \frac{\sqrt 3}{3}.
    \]
    故 $MN$ 与 $AC_1$ 所成角的余弦值为 $\boxed{\dfrac{\sqrt 3}{3}}$。
}

\rmkb{
    \textbf{本章要点}：
    \begin{itemize}
        \item \textbf{求参数}（线段比、点位置）——先写\textbf{参数化坐标}，再让几何条件转化为关于参数的方程；
        \item \textbf{折叠问题}——关键是\textbf{折叠后的坐标}：围绕折叠轴转动，用二面角大小写出旋转矩阵（或简单正交分解）；
        \item \textbf{异面直线角}的计算只需方向向量夹角\textbf{取绝对值再取锐角}，不涉及法向量。
    \end{itemize}
}
